\documentclass[12pt,a4paper]{article}

\usepackage[utf8]{inputenc}
\usepackage[T1]{fontenc}
\usepackage[brazil]{babel}
\usepackage{lmodern}
\usepackage{geometry}
\usepackage{graphicx}
\usepackage{amsmath}
\usepackage{amssymb}
\usepackage{hyperref}

\geometry{margin=2.5cm}

\title{Relatorio de Sistemas Tolerantes a Falhas Distribuidos}
\author{Nome do Autor}
\date{\today}

\begin{document}

\maketitle

\begin{abstract}
Este documento apresenta o relatorio do trabalho desenvolvido na disciplina de
Sistemas Tolerantes a Falhas Distribuidos.
\end{abstract}

\section{Introducao}

Apresente o contexto do problema, a motivacao do trabalho e os objetivos do
relatorio.

\section{Fundamentacao Teorica}

Descreva os conceitos principais relacionados a sistemas distribuidos e
tolerancia a falhas.

\section{Desenvolvimento}

Explique a arquitetura, as decisoes de implementacao e as etapas realizadas no
projeto.

\section{Resultados}

Apresente os resultados obtidos, testes executados e analise do comportamento do
sistema.

\section{Conclusao}

Retome os objetivos do trabalho, resuma os principais aprendizados e indique
possiveis melhorias futuras.

\bibliographystyle{plain}
\bibliography{referencias}

\end{document}
